%==============================================================================
% Apresentacao de slides - Simulador de NoC SPIN (Fat-Tree) em SystemC
% Formato: Beamer (LaTeX)
% Compilacao: pdflatex apresentacao.tex (2x)  ou  Overleaf
%==============================================================================
\documentclass[aspectratio=169]{beamer}

% ----------------------------- Pacotes --------------------------------------
\usepackage[utf8]{inputenc}
\usepackage[T1]{fontenc}
\usepackage[brazil]{babel}
\usepackage{graphicx}
\usepackage{siunitx}
\usepackage{booktabs}
\usepackage{listings}
\usepackage{xcolor}
\usepackage{tikz}
\usetikzlibrary{positioning,arrows.meta,calc,shapes.geometric}

% ----------------------------- Tema -----------------------------------------
\usetheme{Madrid}
\usecolortheme{seahorse}
\setbeamertemplate{navigation symbols}{}
\setbeamertemplate{caption}[numbered]
\setbeamercovered{transparent}

% --------------------- Estilo para logs de simulacao ------------------------
\definecolor{logbg}{RGB}{245,245,245}
\lstdefinestyle{log}{
  basicstyle=\ttfamily\tiny,
  backgroundcolor=\color{logbg},
  breaklines=true,
  columns=fullflexible,
  frame=single,
  framesep=3pt,
  literate={->}{{$\rightarrow$}}1 {>>}{{>>}}2 {<<}{{<<}}2,
}

% ----------------------------- Metadados ------------------------------------
\title[NoC SPIN em SystemC]{Simulador \emph{Cycle-Accurate} de uma Rede-em-Chip
SPIN \\ Topologia Fat-Tree Quaternária em SystemC}
\subtitle{Comutação Wormhole, Controle de Fluxo por Créditos e Arbitragem Probabilística}
\author[Breno, Felipe, Fernando, Gustavo]{Breno Barbosa de Oliveira \and Felipe
Marley de Oliveira Gomes \and Fernando Gomes de Almeida \and Gustavo de Freitas
Rodrigues}
\institute[UFRN -- DIMAp]{Curso de Ciência da Computação \\ Departamento de
Informática e Matemática Aplicada (DIMAp) \\ Universidade Federal do Rio Grande
do Norte}
\date{Trabalho Prático 2 --- Redes-em-Chip (NoC)}

\begin{document}

% ============================== TITULO ======================================
\begin{frame}
  \titlepage
\end{frame}

\begin{frame}{Roteiro}
  \tableofcontents
\end{frame}

% ============================== INTRODUCAO ==================================
\section{Introdução}

\begin{frame}{Introdução --- Por que Redes-em-Chip?}
  \begin{itemize}
    \item SoCs modernos integram dezenas a centenas de núcleos em um mesmo chip.
    \item \textbf{Barramentos} não escalam (banda, energia); \textbf{fios
    dedicados} explodem em área de roteamento.
    \item \alert{Rede-em-Chip (NoC):} comunicação comutada por pacotes, com
    roteadores e enlaces --- como uma rede de computadores em miniatura.
  \end{itemize}
  \vfill
  \begin{block}{Objetivo do trabalho}
    Projetar, implementar e \textbf{verificar por execução} um simulador
    \emph{cycle-accurate} (nível RTL) da NoC \textbf{SPIN}, em \textbf{SystemC},
    com topologia \emph{fat-tree} quaternária.
  \end{block}
\end{frame}

\begin{frame}{Introdução --- Conceitos-chave da SPIN}
  \begin{columns}[T]
    \begin{column}{0.5\textwidth}
      \textbf{Três mecanismos de hardware:}
      \begin{itemize}
        \item \textbf{Wormhole:} mensagem dividida em \emph{flits}
        (\textsc{header}/\textsc{payload}/\textsc{tail}).
        \item \textbf{Créditos:} só transmite se há espaço no \emph{buffer}
        vizinho $\Rightarrow$ sem perda de flit.
        \item \textbf{Árbitro probabilístico:} sorteio a cada ciclo evita
        \emph{starvation}.
      \end{itemize}
    \end{column}
    \begin{column}{0.5\textwidth}
      \textbf{Roteamento em duas fases:}
      \begin{itemize}
        \item \alert{Subida adaptativa:} escolhe o enlace ascendente com mais
        créditos livres.
        \item \alert{Descida determinística:} caminho único dado pelo endereço
        hierárquico do destino.
      \end{itemize}
    \end{column}
  \end{columns}
\end{frame}

% ============================ IMPLEMENTACAO =================================
\section{Implementação e Diagrama de Blocos}

\begin{frame}{Implementação --- Características}
  \begin{columns}[T]
    \begin{column}{0.54\textwidth}
      \begin{itemize}
        \item \textbf{SystemC 3.0.2} (Accellera), C++17, \texttt{g++}~16.
        \item Modelo \emph{cycle-accurate}: processos \texttt{SC\_METHOD}
        sensíveis à borda do relógio (\SI{10}{ns}, \SI{100}{MHz}).
        \item \textbf{8 roteadores} (4 folha + 4 raiz) e \textbf{16 terminais}.
        \item \emph{Buffers} FIFO de profundidade 4 por porta.
        \item Máquina \emph{wormhole} com tabela de travas
        (\texttt{input\_locks}).
      \end{itemize}
    \end{column}
    \begin{column}{0.46\textwidth}
      \footnotesize
      \begin{block}{Módulos}
        \begin{tabular}{@{}ll@{}}
          \texttt{spin\_flit}    & tipo do flit \\
          \texttt{terminal}      & injeta/recebe \\
          \texttt{rspin\_router} & roteador + árbitro \\
          \texttt{spin\_fat\_tree} & topologia (top) \\
          \texttt{main}          & testbench \\
        \end{tabular}
      \end{block}
    \end{column}
  \end{columns}
\end{frame}

\begin{frame}{Diagrama de Blocos --- Topologia Fat-Tree}
  \begin{figure}
  \centering
  \resizebox{0.92\textwidth}{!}{%
  \begin{tikzpicture}[
    every node/.style={font=\scriptsize},
    root/.style={draw,fill=blue!12,minimum width=0.85cm,minimum height=0.6cm,rounded corners=2pt},
    leaf/.style={draw,fill=green!12,minimum width=0.85cm,minimum height=0.6cm,rounded corners=2pt},
    term/.style={draw,fill=orange!18,minimum width=0.5cm,minimum height=0.42cm,font=\tiny,inner sep=1pt},
    ]
    \def\ls{2.4}\def\ts{0.58}
    \foreach \i in {0,1,2,3} { \node[root] (R\i) at (\ls*\i,3.2) {$R_\i$}; }
    \foreach \i in {0,1,2,3} { \node[leaf] (L\i) at (\ls*\i,1.6) {$L_\i$}; }
    \foreach \l in {0,1,2,3} \foreach \r in {0,1,2,3} {
        \draw[gray!55,line width=0.3pt] (L\l.north) -- (R\r.south); }
    \foreach \l in {0,1,2,3} \foreach \p in {0,1,2,3} {
        \pgfmathtruncatemacro{\tid}{4*\l+\p}
        \pgfmathsetmacro{\xx}{\ls*\l + \ts*(\p-1.5)}
        \node[term] (T\tid) at (\xx,0.1) {\tid};
        \draw[gray!70,line width=0.3pt] (L\l.south) -- (T\tid.north); }
    \node[anchor=east] at (-1.5,3.2) {Nível 2 (raiz)};
    \node[anchor=east] at (-1.5,1.6) {Nível 1 (folha)};
    \node[anchor=east] at (-1.5,0.1) {Terminais};
    % destaca o caminho do exemplo: T1 -> L0 -> R0 -> L3 -> T14
    \draw[red,very thick,-{Latex}] (T1.north) -- (L0.south);
    \draw[red,very thick,-{Latex}] (L0.north) -- (R0.south);
    \draw[red,very thick,-{Latex}] (R0.south) -- (L3.north);
    \draw[red,very thick,-{Latex}] (L3.south) -- (T14.north);
  \end{tikzpicture}}
  \end{figure}
  \centering\small Em \textcolor{red}{vermelho}: caminho do exemplo
  $T_1 \to L_0 \to R_0 \to L_3 \to T_{14}$.
\end{frame}

\begin{frame}{Diagrama de Blocos --- Pipeline do Roteador}
  \centering
  \begin{tikzpicture}[node distance=3.2mm,every node/.style={font=\scriptsize},
    stage/.style={draw,fill=blue!8,rounded corners=3pt,minimum width=8.2cm,minimum height=0.6cm,align=left,inner sep=3pt},
    arr/.style={-{Latex[length=2mm]},thick}]
    \node[stage] (s1) {\textbf{1.}~Atualiza créditos recebidos};
    \node[stage,below=of s1] (s2) {\textbf{2.}~Lê flits válidos dos enlaces para os \emph{buffers} FIFO};
    \node[stage,below=of s2,fill=orange!15] (s3) {\textbf{3. Árbitro:} sorteia (\texttt{rand\%100}) a ordem das 8 portas};
    \node[stage,below=of s3,fill=green!15] (s4) {\textbf{4a.}~\emph{Header}: aloca canal (UP adaptativo / DOWN determinístico) e \alert{tranca}};
    \node[stage,below=of s4,fill=green!15] (s5) {\textbf{4b.}~Transmite se há crédito; \emph{Tail} \alert{libera} o canal};
    \node[stage,below=of s5] (s6) {\textbf{5.}~Escreve saídas e devolve créditos};
    \draw[arr] (s1)--(s2);\draw[arr] (s2)--(s3);\draw[arr] (s3)--(s4);
    \draw[arr] (s4)--(s5);\draw[arr] (s5)--(s6);
  \end{tikzpicture}
  \vfill
  \small Executado por \texttt{routing\_logic} a cada borda de relógio.
\end{frame}

% ============================== RESULTADO ===================================
\section{Resultado de Simulação}

\begin{frame}[fragile]{Resultado --- Cenário e Log de Execução Real}
  \small Cenário: \texttt{configure\_traffic(1,14)} --- $T_1$ envia 4 flits a
  $T_{14}$ (sub-árvores distintas). Saída de console \alert{copiada da execução}:
  \vspace{2mm}
\begin{lstlisting}[style=log]
[ PACOTE INICIADO ] @ 30 ns | Origem: 1 Destino: 14
  [ARBITRO] Roteador Nivel 1 ID 0 | caiu em 50% (Prioridade UP) -> Atendendo Entrada 1
      -> [HW] Roteamento Adaptativo trancou porta UP 0
    >> @ 40 ns | Roteador Nivel 1 (ID 0) despachou [HEADER]  -> Porta UP 0
    >> @ 50 ns | Roteador Nivel 2 (ID 0) despachou [HEADER]  -> Porta DOWN 3
    >> @ 60 ns | Roteador Nivel 1 (ID 3) despachou [HEADER]  -> Porta DOWN 2
    <<< @ 70 ns  | Terminal 14 RECEBEU [HEADER]  | Dado: >> INICIO_DA_MENSAGEM <<
    <<< @ 80 ns  | Terminal 14 RECEBEU [PAYLOAD] | Dado: Ola mundo :) #1
    <<< @ 90 ns  | Terminal 14 RECEBEU [PAYLOAD] | Dado: Ola mundo :) #2
    <<< @ 100 ns | Terminal 14 RECEBEU [TAIL]    | Dado: >> FIM_DA_MENSAGEM <<
        -> (Pacote completamente recebido e remontado!)
\end{lstlisting}
\end{frame}

\begin{frame}{Resultado --- Pipeline Sistólico Wormhole}
  \centering
  Ocupação de cada enlace por ciclo (reconstruída do log real):
  \vspace{2mm}

  \renewcommand{\arraystretch}{1.2}
  \begin{tabular}{@{}c|cccc@{}}
    \toprule
    \textbf{t (ns)} & $T_1\!\to\!L_0$ & $L_0\!\to\!R_0$ & $R_0\!\to\!L_3$ & $L_3\!\to\!T_{14}$ \\
    \midrule
    40  & \texttt{PAY\#1} & \texttt{HEADER} & --              & --              \\
    50  & \texttt{PAY\#2} & \texttt{PAY\#1} & \texttt{HEADER} & --              \\
    60  & \texttt{TAIL}   & \texttt{PAY\#2} & \texttt{PAY\#1} & \texttt{HEADER} \\
    70  & --              & \texttt{TAIL}   & \texttt{PAY\#2} & \texttt{PAY\#1} \\
    80  & --              & --              & \texttt{TAIL}   & \texttt{PAY\#2} \\
    90  & --              & --              & --              & \texttt{TAIL}   \\
    \bottomrule
  \end{tabular}

  \vfill
  \small A diagonal evidencia o avanço \emph{wormhole}: 1 salto por ciclo.
  Latência do header = 4 ciclos; pacote completo entregue em \SI{100}{ns}.
\end{frame}

\begin{frame}{Resultado --- Bateria de Casos de Teste (CTest)}
  \centering
  Testes autoverificáveis: cada cenário compara os pacotes remontados com o
  esperado.
  \vspace{2mm}

  \renewcommand{\arraystretch}{1.3}
  \footnotesize
  \begin{tabular}{@{}lll c@{}}
    \toprule
    \textbf{Cenário} & \textbf{Fluxos} & \textbf{Verifica} & \textbf{Status} \\
    \midrule
    Fluxo único    & $T_1\!\to\!T_{14}$        & caminho completo via raiz       & \textcolor{green!55!black}{\textbf{OK}} \\
    Fluxo local    & $T_0\!\to\!T_3$           & roteamento na mesma folha       & \textcolor{green!55!black}{\textbf{OK}} \\
    Convergência   & $T_0,T_8\!\to\!T_7$       & sem perda sob contenção         & \textcolor{green!55!black}{\textbf{OK}} \\
    Permutação     & 8 fluxos disjuntos        & $N$ fluxos sem \emph{deadlock}  & \textcolor{green!55!black}{\textbf{OK}} \\
    Muitos-p/-um   & $T_0..T_3\!\to\!T_{10}$   & serialização justa              & \textcolor{green!55!black}{\textbf{OK}} \\
    \bottomrule
  \end{tabular}

  \vfill
  \alert{Resultado: 5/5 testes aprovados.} A convergência confirma a correção do
  defeito de contenção (perda de cauda quando dois fluxos disputam a saída).
\end{frame}

% ============================== CONCLUSOES ==================================
\section{Conclusões}

\begin{frame}{Conclusões}
  \begin{itemize}
    \item Implementamos e \textbf{executamos} um simulador \emph{cycle-accurate}
    da NoC SPIN (\emph{fat-tree} quaternária) em SystemC.
    \item Validamos os três mecanismos centrais:
    \begin{itemize}
      \item \textbf{Wormhole:} alocação no \emph{header}, fluxo do corpo,
      liberação no \emph{tail};
      \item \textbf{Créditos:} entrega sem perda de flits;
      \item \textbf{Roteamento:} subida adaptativa + descida determinística;
      \item \textbf{Árbitro:} intervenção probabilística a cada salto.
    \end{itemize}
    \item A verificação por execução \alert{expôs e corrigiu um bug real}
    (operador \texttt{==} ignorava o campo \texttt{data}, fazendo o
    \texttt{sc\_signal} suprimir o 2º \emph{payload}).
  \end{itemize}
  \vfill
  \begin{block}{Trabalhos futuros}
    Múltiplos fluxos concorrentes; métricas de latência/vazão; geração de
    \emph{traces} VCD; cargas sintéticas (uniforme, \emph{hotspot}).
  \end{block}
\end{frame}

\begin{frame}[plain]
  \centering
  \vfill
  {\Huge Obrigado!}\\[4mm]
  {\large Perguntas?}
  \vfill
  \small Breno \quad Felipe \quad Fernando \quad Gustavo --- UFRN/DIMAp
  \vfill
\end{frame}

\end{document}
